% ========================================================
% INVESTIGACIÓN BIBLIOGRÁFICA: 3 TEMAS, 3 PAPERS POR INTEGRANTE
% ========================================================
\section{Investigación bibliográfica}

\subsection{Pregunta guía y criterios de selección}

La revisión parte de una pregunta que el experimento de PC2 podrá contrastar: \textit{¿cómo afectan
la granularidad del trabajo y la estrategia de reducción al rendimiento y a la equivalencia numérica
de K-means sobre más de un millón de registros?}

Con esa pregunta como filtro, seleccionamos artículos publicados desde 2021 cuyo tema central fuera
un modelo de machine learning acelerado con concurrencia. Todos provienen de revistas o conferencias
arbitradas de IEEE, ACM, Springer, Elsevier, MDPI o PMLR. Verificamos el DOI, los autores y el año
contra el registro de cada editor, y preferimos versiones de acceso abierto para que cualquier lector
pueda consultar la fuente.
Seis de los nueve venues están en el primer cuartil de SCImago o en el rango A de ICORE
(Tabla~\ref{tab:ranking}); las excepciones son la revista \textit{Algorithms} (Q2) y las conferencias
ICPP y PPoPP (rango B).
Los trabajos clásicos de la línea, anteriores a esa ventana, se citan como antecedentes en la
sección~\ref{sec:antecedentes}.

% Tabla: posición de los venues (cuartil SCImago para revistas, rango ICORE para conferencias).
\begin{table}[H]
\caption{Posición de los venues de los nueve artículos}\label{tab:ranking}
\small
\begin{tabularx}{\textwidth}{@{}L{3.6cm}YL{3.4cm}@{}}
\toprule
\textbf{Artículo} & \textbf{Venue} & \textbf{Indicador} \\
\midrule
\textcite{nigro2022}      & \textit{Algorithms} (MDPI), revista                       & SJR Q2 \\
\textcite{kwedlo2021}     & \textit{IEEE Access} (IEEE), revista                      & SJR Q1 \\
\textcite{kwedlo2024}     & ICPP (ACM), conferencia                                   & ICORE B \\
\addlinespace
\textcite{hess2023}       & AISTATS (PMLR), conferencia                               & ICORE A \\
\textcite{bellavita2026}  & IPDPS (IEEE), conferencia                                 & ICORE A \\
\textcite{xu2021}         & \textit{Computer Networks} (Elsevier), revista            & SJR Q1 \\
\addlinespace
\textcite{li2023}         & \textit{Data Mining and Knowledge Discovery} (Springer), revista & SJR Q1 \\
\textcite{bellavita2025}  & PPoPP (ACM), conferencia                                  & ICORE B \\
\textcite{voevodski2024}  & \textit{ACM Transactions on Knowledge Discovery from Data}, revista & SJR Q1 \\
\bottomrule
\end{tabularx}

{\footnotesize\textit{Nota.} Revistas: mejor cuartil del SCImago Journal Rank 2024, calculado con
datos de Scopus. Conferencias: rango de ICORE 2026 (A* es el más alto, seguido de A, B y C). PPoPP
figuraba como A en CORE 2021 y pasó a B en CORE 2023 porque no presentó datos para su revisión.
Consultado el 13 de septiembre de 2026. Los grupos de filas corresponden a los Temas 1, 2 y 3.\par}
\end{table}


La Tabla~\ref{tab:reparto} muestra el reparto de artículos. En cada ficha, el campo \emph{aporte de
la concurrencia} recoge el resultado que el artículo atribuye al paralelismo; su relación con
nuestro diseño se discute al cierre de cada tema.

\begin{table}[H]
\caption{Reparto de los nueve artículos entre temas e integrantes}\label{tab:reparto}
\small
\begin{tabularx}{\textwidth}{@{}L{4.4cm}L{3.4cm}Y@{}}
\toprule
\textbf{Tema} & \textbf{Responsable} & \textbf{Artículos} \\
\midrule
1. Paralelismo en memoria compartida & \nombreTercero
  & \textcite{nigro2022}; \textcite{kwedlo2021}; \textcite{kwedlo2024} \\
\addlinespace
2. Sincronización distribuida & \nombreDayana
  & \textcite{hess2023}; \textcite{bellavita2026}; \textcite{xu2021} \\
\addlinespace
3. Poda, localidad y aproximación & \nombreRody
  & \textcite{li2023}; \textcite{bellavita2025}; \textcite{voevodski2024} \\
\bottomrule
\end{tabularx}
\end{table}

% --------------------------------------------------------
\subsection{Tema 1: paralelismo de datos y sincronización en memoria compartida}
% --------------------------------------------------------

\textbf{Responsable:} \nombreTercero.

Es el tema más cercano a nuestra implementación: una máquina, varios núcleos y el dataset completo
en memoria. Interesa cómo se reparte la asignación de puntos entre hilos, cómo se combinan los
resultados parciales y cuánto pesa la coordinación frente al cálculo.

\tarjeta{nigro2022}{K-means paralelo en Java}{
  titulo={Performance of Parallel K-Means Algorithms in Java},
  autores={Libero Nigro},
  fuente={2022, \textit{Algorithms} (MDPI), 15(4), 117},
  link={\url{https://doi.org/10.3390/a15040117}},
  motivacion={K-means se encarece al crecer los registros, la dimensión y el número de clusters,
    mientras que las máquinas multinúcleo ya son hardware común.},
  problema={Paralelizar K-means en una sola máquina y determinar qué modelo de concurrencia de Java
    rinde más sin alterar el resultado.},
  propuesta={Compara dos versiones de Lloyd: map/reduce con \textit{parallel streams}, donde el
    runtime decide el particionado, y actores sin locks con Theatre, donde cada hilo es una unidad
    gruesa asignada a un núcleo.},
  algoritmo={\paso{1}Inicializar $K$ centroides.
    \paso{2}Asignar cada punto en paralelo (\textit{map}).
    \paso{3}Sumar puntos y conteos por cluster (\textit{reduce}).
    \paso{4}Recalcular centroides hasta que se muevan menos de $10^{-10}$ o se llegue a 1000
    iteraciones.},
  stack={Java, Streams API y biblioteca de actores Theatre; patrones map/reduce y modelo de
    actores.},
  aporte={Con US Census 1990 (2.458.285 registros, 68 dimensiones) y 16 hilos, logra un speedup
    cercano a 11 con \textit{streams} ($K=80$) y de 10,63 con Theatre ($K=120$). Ambas versiones
    reproducen exactamente los centroides y el número de iteraciones de la versión secuencial.},
  integrante={\nombreTercero},
}

\tarjeta{kwedlo2021}{K-means acelerado en memoria compartida}{
  titulo={Accelerated K-Means Algorithms for Low-Dimensional Data on Parallel Shared-Memory
    Systems},
  autores={Wojciech Kwedlo, Michał Łubowicz},
  fuente={2021, \textit{IEEE Access}, 9, 74286--74301},
  link={\url{https://doi.org/10.1109/ACCESS.2021.3080821}},
  motivacion={Los K-means acelerados evitan distancias innecesarias, pero cuestan más de
    paralelizar: la carga de trabajo por punto deja de ser uniforme.},
  problema={Paralelizar un K-means exacto para datos de baja dimensión sin que la gestión de
    tareas limite la escalabilidad.},
  propuesta={Versión paralela del \textit{filtering algorithm}. Un kd-tree descarta los centroides
    que no pueden ser los más cercanos a ningún punto de una celda; la construcción del árbol y su
    recorrido se reparten en tareas OpenMP, una por llamada recursiva.},
  algoritmo={\texttt{Filter}($u$, $C$):
    \paso{1}En una hoja, iterar Lloyd solo con los candidatos $C$.
    \paso{2}En un nodo interno, quitar de $C$ los centroides más lejanos de toda la celda que el más
    cercano a su punto medio.
    \paso{3}Con un solo candidato, asignar la celda entera usando la suma y el conteo del nodo.
    \paso{4}Si quedan varios, recorrer los dos hijos como tareas paralelas.},
  stack={C/C++ con OpenMP \textit{tasking}; divide y vencerás sobre memoria compartida.},
  aporte={En 24 núcleos, con 6 datasets sintéticos (hasta $10^9$ puntos) y 7 reales, alcanza una
    eficiencia paralela satisfactoria o alta y tiempos mucho menores que Lloyd paralelo, una versión
    en GPU y dos métodos basados en la desigualdad triangular. La ventaja cae rápido al aumentar la
    dimensión.},
  integrante={\nombreTercero},
}

\tarjeta{kwedlo2024}{IMM paralelo en memoria compartida}{
  titulo={Parallel Iterative Mistake Minimization (IMM) clustering algorithm for shared-memory
    systems},
  autores={Wojciech Kwedlo},
  fuente={2024, ICPP '24 (ACM), 1--10},
  link={\url{https://doi.org/10.1145/3673038.3673057}},
  motivacion={Un árbol de decisión compacto vuelve legibles los clusters de K-means, pero
    construirlo con IMM es costoso.},
  problema={Paralelizar IMM, el algoritmo que genera ese árbol explicativo a partir de los
    centroides.},
  propuesta={Dos versiones con OpenMP. Una paraleliza el bucle externo sobre las features, como la
    implementación de referencia; la otra agrega paralelismo anidado para buscar en paralelo el
    corte óptimo dentro de cada feature.},
  algoritmo={\paso{1}Obtener los centroides con K-means.
    \paso{2}En cada nodo, evaluar los umbrales de cada feature y contar los puntos que quedarían
    separados de su centroide.
    \paso{3}Dividir por el corte con menos errores.
    \paso{4}Repetir hasta que cada hoja tenga un solo centroide.},
  stack={C/C++ con OpenMP y paralelismo anidado; nodos de 48 núcleos y 9 datasets sintéticos.},
  aporte={La versión con paralelismo anidado es entre 1,3 y 15 veces más rápida que la
    re-implementación de referencia, según la dimensión del espacio de features.},
  integrante={\nombreTercero},
}

\paragraph{Síntesis.} Los tres trabajos corren sobre memoria compartida, pero reparten unidades
de trabajo muy distintas (Tabla~\ref{tab:sintesis1}). Esa elección pesa tanto como el número de
hilos: el mismo algoritmo IMM cambia hasta 15 veces de rendimiento según el nivel en que se abre el
paralelismo \parencite{kwedlo2024}. Para PC2, esto significa variar el tamaño de bloque además del
número de goroutines. De \textcite{nigro2022} tomamos el criterio de corrección: la versión
concurrente debe reproducir los centroides de la secuencial.

\begin{table}[H]
\caption{Comparación de los artículos del Tema 1}\label{tab:sintesis1}
\small
\begin{tabularx}{\textwidth}{@{}L{3.3cm}YYY@{}}
\toprule
\textbf{Artículo} & \textbf{Mecanismo} & \textbf{Unidad que se reparte} & \textbf{Aceleración reportada} \\
\midrule
\textcite{nigro2022} & \textit{Parallel streams} y actores & Bloques regulares de puntos
  & Speedup cercano a 11 con 16 hilos, frente a la versión secuencial \\
\addlinespace
\textcite{kwedlo2021} & Tareas OpenMP recursivas & Subárboles del kd-tree, de tamaño variable
  & Menor tiempo que Lloyd paralelo, GPU y dos métodos con desigualdad triangular \\
\addlinespace
\textcite{kwedlo2024} & OpenMP anidado & Features y umbrales candidatos
  & Entre 1,3 y 15 veces frente a la re-implementación de referencia \\
\bottomrule
\end{tabularx}
\end{table}

% --------------------------------------------------------
\subsection{Tema 2: sincronización en entornos distribuidos}
% --------------------------------------------------------

\textbf{Responsable:} \nombreDayana.

Cuando el trabajo se reparte entre máquinas, el costo dominante deja de ser el cálculo. Cada
barrera de sincronización obliga a comunicar resultados y a esperar al participante más lento, y
los tres artículos de este tema buscan pagar menos por ella.

\tarjeta{hess2023}{K-means distribuido en pocas rondas}{
  titulo={Fast Distributed $k$-Means with a Small Number of Rounds},
  autores={Tom Hess, Ron Visbord, Sivan Sabato},
  fuente={2023, AISTATS, PMLR 206, 850--874},
  link={\url{https://proceedings.mlr.press/v206/hess23a.html}},
  motivacion={En K-means distribuido, una ronda de comunicación entre máquinas cuesta mucho más que
    el cómputo local.},
  problema={Obtener un clustering con garantía de aproximación usando pocas rondas de
    comunicación.},
  propuesta={SOCCER, un esquema coordinador--workers cuyo factor de aproximación y número de rondas
    dependen solo de la capacidad del coordinador. Un criterio de parada incorporado ahorra rondas
    cuando los datos lo permiten.},
  algoritmo={\paso{1}Cada máquina envía una muestra de sus puntos al coordinador.
    \paso{2}El coordinador agrupa la muestra y obtiene centros.
    \paso{3}Las máquinas descartan los puntos ya bien cubiertos por esos centros.
    \paso{4}Se repite con los restantes hasta la parada y se combinan los centros.},
  stack={Coordinador--workers (maestro--esclavo); código oficial en Python
    (\url{https://github.com/selotape/distributed_k_means}); comparación con k-means$\|$
    \parencite{bahmani2012}.},
  aporte={En muchos datasets naturales bastan de 1 a 4 rondas. El costo del clustering suele ser
    mejor que el de k-means$\|$, aun cuando a este se le permiten más rondas, y el tiempo de cómputo
    en las máquinas es bastante menor.},
  integrante={\nombreDayana},
  opinion={Al analizar este trabajo, considero que el principal aporte de la programación
    concurrente no está en acelerar únicamente el cálculo de K-means, sino en reducir el costo de
    coordinación entre máquinas. En sistemas distribuidos, cada ronda de comunicación requiere
    sincronización entre participantes, generando tiempos de espera que pueden convertirse en el
    principal cuello de botella. SOCCER aborda este problema reduciendo la cantidad de rondas
    necesarias mediante un esquema coordinador--workers y un mecanismo de parada adaptativa,
    logrando que en muchos casos naturales sean suficientes pocas rondas de comunicación.
    \par El beneficio más importante es que permite escalar K-means cuando los datos ya están
    distribuidos entre varias máquinas, aprovechando el procesamiento paralelo sin depender de una
    comunicación constante. Sin embargo, esta estrategia también introduce limitaciones: el
    coordinador concentra una parte importante del procesamiento y su capacidad de cómputo y memoria
    influye directamente en el rendimiento del algoritmo. Además, aunque se reduzcan las
    comunicaciones, no desaparece el costo asociado a transmitir muestras y sincronizar resultados.
    \par Considero que este trabajo demuestra que la concurrencia no consiste simplemente en añadir
    más máquinas, sino en diseñar algoritmos que minimicen la necesidad de coordinación. En
    aplicaciones reales, sería necesario evaluar si el costo de comunicación justifica la
    distribución, ya que en datasets pequeños el paralelismo podría incluso añadir complejidad
    innecesaria.},
}

\tarjeta{bellavita2026}{Kernel K-means con poca comunicación entre GPUs}{
  titulo={Communication-Avoiding Linear Algebraic Kernel K-Means on GPUs},
  autores={Julian Bellavita, Matthew Rubino, Nakul Iyer, Andrew Chang, Aditya Devarakonda, Flavio
    Vella, Giulia Guidi},
  fuente={2026, IEEE IPDPS; se cita la versión de arXiv mientras no haya DOI},
  link={\url{https://arxiv.org/abs/2601.17136}},
  motivacion={La versión para una sola GPU \parencite{bellavita2025} no pasa de unas 80.000
    muestras, porque la matriz kernel no cabe en memoria.},
  problema={Escalar Kernel K-means a millones de puntos en muchas GPUs sin que la comunicación
    domine el tiempo.},
  propuesta={Algoritmos de memoria distribuida con particionado 1D y 1.5D, que componen primitivas
    de álgebra lineal comunicando lo mínimo entre GPUs.},
  algoritmo={\paso{1}Particionar la matriz kernel entre GPUs (1D o 1.5D).
    \paso{2}Multiplicar en cada GPU su bloque por la matriz dispersa de pertenencia a clusters.
    \paso{3}Intercambiar solo los resultados parciales necesarios para las distancias.
    \paso{4}Reasignar los puntos y repetir hasta converger.},
  stack={Multi-GPU con memoria distribuida, software Vivaldi
    (\url{https://github.com/CornellHPC/Vivaldi}); particionado de datos y álgebra lineal
    dispersa.},
  aporte={En 256 GPUs alcanza 79,7\,\% de eficiencia media en weak scaling y un speedup medio de 4,2
    en strong scaling. El esquema 1.5D llega a ser 3,6 veces más rápido que el 1D, y baja de más de
    una hora a menos de dos segundos frente a una GPU con ventana deslizante.},
  integrante={\nombreDayana},
  opinion={Al analizar este trabajo, considero que el principal aporte de la programación
    concurrente no está únicamente en distribuir el cálculo de Kernel K-means entre varias GPUs,
    sino en diseñar una arquitectura donde la comunicación entre ellas no limite la escalabilidad.
    Inicialmente podría pensarse que aumentar la cantidad de GPUs siempre produce una mejora
    proporcional; sin embargo, en algoritmos distribuidos el intercambio de matrices y la
    sincronización pueden convertirse en el nuevo cuello de botella. Los autores muestran que
    Kernel K-means presenta altos costos computacionales y de memoria debido al crecimiento
    cuadrático de la matriz kernel, lo que limita las implementaciones en una sola GPU.
    \par El beneficio principal de este trabajo es que la concurrencia permite ejecutar clustering
    sobre datasets de millones de muestras mediante algoritmos distribuidos como el enfoque 1.5D,
    que reduce la comunicación mediante esquemas de particionamiento especializados. Sin embargo,
    esta solución también aumenta la complejidad del sistema, ya que requiere administrar
    distribución de datos, balance de carga y sincronización entre dispositivos. Además, una mala
    estrategia de particionamiento podría generar más comunicación que beneficio.
    \par Considero que este trabajo demuestra que en computación paralela no basta con agregar más
    hardware; es necesario diseñar algoritmos conscientes de la comunicación. La concurrencia aporta
    escalabilidad, pero solo cuando el costo de coordinar múltiples dispositivos es menor que la
    ganancia obtenida por paralelizar el cálculo.},
}

\tarjeta{xu2021}{workers de respaldo dinámicos}{
  titulo={Dynamic backup workers for parallel machine learning},
  autores={Chuan Xu, Giovanni Neglia, Nicola Sebastianelli},
  fuente={2021, \textit{Computer Networks} (Elsevier), 188, 107846},
  link={\url{https://doi.org/10.1016/j.comnet.2021.107846}; versión abierta:
    \url{https://arxiv.org/abs/2004.14696}},
  motivacion={Con un servidor de parámetros síncrono, basta un worker lento (\textit{straggler}) para
    alargar cada iteración. La lentitud puede venir de contención de recursos, actividad del sistema
    operativo o recolección de basura.},
  problema={Esperar solo las $n-b$ actualizaciones más rápidas mitiga el problema, pero el $b$
    óptimo depende del clúster, la carga, el tamaño de lote y la etapa del entrenamiento. Fijarlo de
    antemano exige experimentos costosos y resulta frágil.},
  propuesta={DBW (\textit{Dynamic Backup Workers}) decide en cada iteración cuántas actualizaciones
    esperar, sin conocer el clúster, maximizando la reducción de la pérdida por unidad de tiempo.},
  algoritmo={\paso{1}El servidor envía los parámetros a los $n$ workers.
    \paso{2}Cada worker calcula el gradiente de un mini-lote.
    \paso{3}DBW estima, para cada $k$, la duración de la iteración y la bajada de la pérdida, y
    elige el $k$ con mayor progreso por unidad de tiempo.
    \paso{4}El servidor agrega las primeras $k$ actualizaciones y descarta las demás.},
  stack={PyTorch con backend MPI; servidor de parámetros (coordinador--workers) con barrera de
    sincronización parcial. Experimentos en un clúster real de CPU y GPU (GeForce GTX 1080 Ti, GTX
    Titan X y Tesla V100) con 16 workers: una red convolucional sobre MNIST y ResNet18 sobre CIFAR10.
    Código: \url{https://gitlab.inria.fr/chxu/dbw}.},
  aporte={Elimina el ajuste previo de $b$ y entrena hasta 3 veces más rápido que la mejor
    configuración estática.},
  integrante={\nombreDayana},
  opinion={Al analizar este trabajo, considero que el principal aporte de la programación
    concurrente está en reducir el impacto de los retrasos dentro de los sistemas distribuidos de
    aprendizaje automático. En un esquema tradicional basado en \textit{parameter server}, cada
    iteración depende de recibir las actualizaciones de todos los trabajadores, por lo que un solo
    nodo lento puede retrasar todo el entrenamiento.
    \par La propuesta Dynamic Backup Workers (DBW) muestra que el paralelismo no consiste únicamente
    en agregar más trabajadores, sino en gestionar mejor la sincronización entre ellos. Al permitir
    que el servidor utilice las actualizaciones de los trabajadores más rápidos y considere al resto
    como \textit{backup workers}, se reduce el tiempo perdido esperando nodos lentos. Sin embargo,
    considero que esta ventaja tiene un costo: disminuir la cantidad de trabajadores sincronizados
    puede introducir mayor variabilidad en los gradientes y afectar la estabilidad del entrenamiento.
    \par El aspecto más interesante del trabajo es que convierte un problema de sincronización en una
    decisión adaptativa del sistema. Aun así, no asumiría que DBW siempre es superior: en escenarios
    con baja variabilidad entre trabajadores, la complejidad adicional del mecanismo dinámico podría
    no justificar la mejora obtenida. Evaluaría primero cuánto tiempo se pierde realmente por
    \textit{stragglers} antes de incorporar esta estrategia, ya que el beneficio depende directamente
    del comportamiento del entorno distribuido.},
}

\paragraph{Síntesis.} Cada artículo reduce un componente distinto del costo de sincronizar
(Tabla~\ref{tab:sintesis2}), y no todos conservan el algoritmo original. SOCCER entrega una
aproximación en lugar de las iteraciones de Lloyd, y DBW actualiza los parámetros con información
incompleta. Nuestra barrera con \texttt{sync.WaitGroup} espera a todos los workers, así que el
bloque más lento fija la duración de cada iteración. En un K-means exacto preferimos pagar esa
espera, y acotarla con bloques de tamaño parecido, antes que perder la equivalencia con la versión
secuencial.

\begin{table}[H]
\caption{Comparación de los artículos del Tema 2}\label{tab:sintesis2}
\small
\begin{tabularx}{\textwidth}{@{}L{3.3cm}YYY@{}}
\toprule
\textbf{Artículo} & \textbf{Qué reduce} & \textbf{Qué cede a cambio} & \textbf{Aceleración reportada} \\
\midrule
\textcite{hess2023} & Número de rondas de comunicación & Aproxima; no replica Lloyd
  & De 1 a 4 rondas, con mejor costo que k-means$\|$ \\
\addlinespace
\textcite{bellavita2026} & Volumen comunicado entre GPUs & Particionado más complejo (1D, 1.5D)
  & Speedup medio de 4,2 en strong scaling con 256 GPUs \\
\addlinespace
\textcite{xu2021} & Espera en cada barrera & Descarta las actualizaciones más lentas
  & Hasta 3 veces frente al mejor $b$ estático \\
\bottomrule
\end{tabularx}
\end{table}

% --------------------------------------------------------
\subsection{Tema 3: poda, localidad y aproximación}
% --------------------------------------------------------

\textbf{Responsable:} \nombreRody.

Paralelizar no es la única forma de acelerar K-means. Los artículos de este tema reducen el
trabajo en sí: omiten distancias y transferencias, cambian la representación del cálculo o aceptan
una aproximación con garantía demostrada.

\tarjeta{li2023}{K-means a gran escala en GPU}{
  titulo={Large scale K-means clustering using GPUs},
  autores={Mi Li, Eibe Frank, Bernhard Pfahringer},
  fuente={2023, \textit{Data Mining and Knowledge Discovery} (Springer), 37(1), 67--109},
  link={\url{https://doi.org/10.1007/s10618-022-00869-6}},
  motivacion={La mayoría de las implementaciones de K-means en GPU exigen cargar el dataset completo
    en la memoria de la GPU.},
  problema={Ejecutar K-means exacto sobre datasets mayores que esa memoria sin que las
    transferencias entre CPU y GPU dominen el tiempo.},
  propuesta={ASB K-means (\textit{Asynchronous Selective Batched}) combina las cotas de Hamerly con
    lotes y GPU. Cada punto guarda una cota superior de la distancia a su centro y una inferior de
    la distancia a los demás; por la desigualdad triangular, si la superior no supera a la inferior,
    el punto no puede cambiar de cluster y se omite. Solo los puntos restantes viajan a la GPU, por
    lotes, mientras la CPU prepara el lote siguiente.},
  algoritmo={\paso{1}Inicializar los centros con k-means++ y las cotas de cada punto (CPU).
    \paso{2}Armar un lote con los puntos cuyas cotas no descartan un cambio de cluster (CPU).
    \paso{3}Calcular en paralelo la distancia de cada punto del lote a todos los centros, y guardar
    el más cercano y la distancia al segundo (GPU).
    \paso{4}Acumular sumas y conteos por cluster y repetir 2--3 hasta cubrir los puntos (CPU).
    \paso{5}Recalcular los centros (CPU), actualizar las cotas con lo que se movió cada uno (GPU) y
    repetir hasta converger.},
  stack={CUDA 10.2 con \textit{streams} asíncronos y memoria \textit{pinned}; GPU NVIDIA Titan X
    (3.584 núcleos, 12~GB) y CPU Intel Core i7 con 64~GB de RAM. Lotes de 35.840 puntos, 10 veces
    los núcleos de la GPU: con lotes mayores, seleccionar en la CPU tarda más que calcular en la GPU,
    que queda esperando. Compara contra RAPIDS cuML, Lloyd en GPU y Elkan y Hamerly en CPU, con los
    mismos centros iniciales y 300 iteraciones.},
  aporte={Produce el mismo resultado que K-means estándar, y tras las primeras iteraciones el
    90\,\% o más de los puntos no cambia de cluster, así que a la GPU llega poco. Es hasta 15 veces
    más rápido que Lloyd en GPU en su mejor caso, pero en las pruebas sintéticas y de imágenes que
    detalla la ventaja es de 1,7 a 2,1 veces. Frente a RAPIDS va de 1,65 veces (vectores de 50
    dimensiones) a 1,09 (200 dimensiones), y en los datasets chicos Elkan con 8 hilos de CPU supera
    ligeramente al Lloyd en GPU. Con los vectores de Twitter ($\approx$1~GB) usa 195~MB de memoria de
    GPU, frente a 1.055~MB de Lloyd en GPU y 4.719~MB de RAPIDS.},
  integrante={\nombreRody},
  opinion={Al principio atribuí el 15x a la GPU, pero la comparación que importa es contra otro
    K-means que también corre en GPU, y ahí la ventaja baja a 1,7--2,1 veces. El beneficio principal
    no está en el paralelismo, sino en calcular menos: la desigualdad triangular descarta los puntos
    que no pueden cambiar de cluster. La concurrencia aporta en otro punto, porque la CPU prepara el
    siguiente lote mientras la GPU calcula, en vez de quedarse ociosa. Veo tres perjuicios. Esa
    coordinación tiene un costo, y con lotes grandes la GPU termina esperando a la CPU. En las
    primeras iteraciones casi todos los puntos cambian de cluster, así que la poda no ahorra nada.
    Y la ventaja se diluye al subir la dimensión; en datasets chicos, Elkan con 8 hilos de CPU le
    gana al Lloyd en GPU, lo que muestra que más hardware paralelo no garantiza más velocidad. Para
    nuestro caso, las 6 dimensiones favorecen la poda y las cotas de cada punto caben en los
    acumuladores privados de cada worker (unos 45~MB para 2,8 millones de viajes). Aun así, primero
    mediría el Lloyd concurrente sin poda, para no mezclar la mejora algorítmica con el speedup de la
    concurrencia.},
}

\tarjeta{bellavita2025}{Popcorn, Kernel K-means en GPU}{
  titulo={Popcorn: Accelerating Kernel K-means on GPUs through Sparse Linear Algebra},
  autores={Julian Bellavita, Thomas Pasquali, Laura Del Rio Martin, Flavio Vella, Giulia Guidi},
  fuente={2025, PPoPP '25 (ACM), 426--440},
  link={\url{https://doi.org/10.1145/3710848.3710887}; versión abierta:
    \url{https://arxiv.org/abs/2501.05587}},
  motivacion={Kernel K-means encuentra clusters que no son linealmente separables, pero su costo
    crece de forma cuadrática con el tamaño del dataset.},
  problema={Acelerar Kernel K-means en GPU con poco esfuerzo de programación.},
  propuesta={Reformular el algoritmo con primitivas de álgebra lineal dispersa (SpMM y SpMV). Es la
    primera implementación de código abierto de Kernel K-means en GPU.},
  algoritmo={\paso{1}Calcular la matriz kernel $\mathbf{K}$.
    \paso{2}Codificar la asignación en una matriz dispersa indicadora $\mathbf{V}$.
    \paso{3}Obtener las distancias con el producto disperso $\mathbf{K}\mathbf{V}$.
    \paso{4}Actualizar $\mathbf{V}$ y repetir.},
  stack={CUDA 12.2 con cuBLAS (matriz kernel, eligiendo GEMM o SYRK según n/d), cuSPARSE (SpMM y
    SpMV) y thrust; menos de 50 líneas de CUDA escritas a mano. GPU NVIDIA A100 con CPU AMD EPYC
    7763 de 64 núcleos; 6 datasets de libSVM de hasta 78.823 puntos, 30 iteraciones fijas y kernel
    polinomial. Artefacto reproducible publicado en Zenodo.},
  aporte={Hasta 123,8 veces más rápido que PRMLT, una implementación en MATLAB de un solo hilo.
    Frente a una implementación CUDA propia sin matrices dispersas, es de 1,6 a 2,6 veces más rápido
    en el algoritmo completo; la base CUDA, a su vez, supera a PRMLT normalmente entre 11 y 21 veces.},
  integrante={\nombreRody},
  opinion={El 123,8x me parece parcialmente engañoso: compara una A100 contra MATLAB de un solo hilo,
    en una máquina con 64 núcleos. La comparación honesta es contra otra GPU, y ahí la ventaja es de
    1,6 a 2,6 veces. El \textit{kernel trick} es muy potente, porque compara los puntos como si
    vivieran en un espacio de muchas más dimensiones sin construirlo, pero lo cobra en memoria: la
    matriz kernel crece con $n^2$. Con 78.823 puntos ya ocupa casi 25~GB en precisión simple, y con
    nuestros 2,8 millones de viajes serían unos 32~TB. Ninguna cantidad de hilos ni de GPU reduce
    ese costo cuadrático: el paralelismo acelera el cálculo, pero no achica lo que hay que guardar.
    Me convence delegar la concurrencia en cuBLAS y cuSPARSE, porque se escribe muy poco CUDA y se
    hereda el trabajo de optimización de NVIDIA. El costo es perder control: el propio artículo
    muestra que SpMM no usa la memoria compartida de la GPU y mueve más datos, y además quedas atado
    a un proveedor. Me gustaría medirlo con submuestras de nuestros viajes: en Go, con
    goroutines, sobre los 15~GB de RAM de mi laptop, donde según mi cálculo la matriz deja de caber
    cerca de los 60.000 puntos, y con una prueba rápida en una GPU RTX 4060 Ti de 16~GB. La idea es mostrar que
    más goroutines bajan el tiempo pero no la memoria, y que la salida pasa por repartir la matriz
    entre varias GPU, como proponen \textcite{bellavita2026}.},
}

\tarjeta{voevodski2024}{clustering $k$ a gran escala con aproximación}{
  titulo={Large-Scale K-Clustering},
  autores={Konstantin Voevodski},
  fuente={2024, \textit{ACM Transactions on Knowledge Discovery from Data}, 18(9), art. 214, 1--23},
  link={\url{https://doi.org/10.1145/3674508}},
  motivacion={Con colecciones de miles de millones de puntos, los algoritmos de clustering con
    garantías teóricas dejan de escalar.},
  problema={Diseñar algoritmos paralelos de $k$-median, $k$-medoids y $k$-means que conserven una
    garantía de calidad frente a la solución óptima.},
  propuesta={Agrupar solo una muestra de semillas y estimar el costo sobre todo el dataset con un
    resumen de la celda de Voronoi de cada semilla. El autor demuestra una aproximación de factor
    constante en instancias \emph{estables}, cuyo clustering óptimo no cambia ante pequeñas
    perturbaciones.},
  algoritmo={\paso{1}Muestrear $t$ semillas al azar.
    \paso{2}Calcular en paralelo la descomposición de Voronoi: cada máquina recibe las semillas y
    devuelve el tamaño y el costo de cada celda, que luego se suman.
    \paso{3}Construir un árbol jerárquico de las semillas con distancia minimax.
    \paso{4}Elegir con programación dinámica la mejor poda de $k$ grupos.},
  stack={Particionado de datos con reducción por suma (parte paralela) y cálculo en memoria de una sola
    máquina (árbol y programación dinámica, $O(t^3)$). El artículo no nombra lenguaje ni framework:
    mide en una CPU Intel Xeon de 2,20~GHz sobre Covertype, Quantum Physics y Protein Homology, y
    proyecta la parte paralela a $m$ máquinas.},
  aporte={Además de la garantía teórica, sus algoritmos de $k$-median y $k$-means encuentran mejores
    clusterings que las construcciones de coresets del estado del arte con muestras del mismo
    tamaño.},
  integrante={\nombreRody},
  opinion={Me parece interesante que cada celda de Voronoi se resuma con conteos y sumas calculados en
    paralelo: es el mismo patrón de nuestros acumuladores, pero con una sola ronda de comunicación en
    lugar de una barrera por iteración. Aun así, no aceptaría una aproximación de hasta 10 veces el
    óptimo cuando Lloyd exacto cabe en mi laptop: las 6 features de 2,8 millones de viajes ocupan
    unos 136~MB. Tampoco creo que nuestros viajes cumplan la estabilidad que exige la garantía. La
    movilidad cambia con feriados, eventos, obras, fines de semana y políticas públicas como el cobro
    por congestión, y esa mezcla produce tipos de viaje que se solapan, justo lo contrario de los
    clusters bien separados que supone el artículo. En cuanto a la concurrencia, el beneficio es
    repartir la pasada $O(n)$ entre muchas máquinas; el perjuicio es que el árbol y la programación
    dinámica, $O(t^3)$, corren en una sola y, por la ley de Amdahl, ponen techo al speedup. Lo que el
    autor afirma es escalamiento débil (más datos en el mismo tiempo), y lo proyecta sin medirlo en
    paralelo; él mismo reconoce que sus algoritmos no son necesariamente más rápidos que las
    alternativas.},
}

\paragraph{Síntesis.} Las cifras de este tema no son comparables entre sí, porque cada artículo
mide algo distinto (Tabla~\ref{tab:sintesis3}). ASB acelera sin perder exactitud, y frente a otro
K-means que también corre en GPU su ventaja es de 1,7 a 2,1 veces en las pruebas que detalla (15 en
el mejor caso): buena parte de la aceleración viene de omitir puntos, no de la GPU. Voevodski cambia
exactitud por una garantía de aproximación y reporta calidad del clustering, no tiempo. Popcorn deja
la lección de evaluación que más nos afecta: 123,8 y 2,6 describen el mismo sistema contra dos
líneas base distintas. Por eso, nuestra tabla de speedup declarará siempre contra qué versión se
mide.

\begin{table}[H]
\caption{Comparación de los artículos del Tema 3}\label{tab:sintesis3}
\small
\begin{tabularx}{\textwidth}{@{}L{3.3cm}YYY@{}}
\toprule
\textbf{Artículo} & \textbf{Qué evita} & \textbf{Resultado frente al exacto} & \textbf{Aceleración reportada} \\
\midrule
\textcite{li2023} & Distancias y transferencias CPU--GPU & Exacto
  & Hasta 15 veces frente a Lloyd en GPU; de 1,09 a 1,65 frente a RAPIDS \\
\addlinespace
\textcite{bellavita2025} & Operaciones densas sobre la matriz kernel & Mismo modelo, otra representación
  & 123,8 veces frente a CPU; 2,6 frente a GPU densa \\
\addlinespace
\textcite{voevodski2024} & Agrupar todos los puntos: agrupa solo una muestra & Aproximación de factor constante
  & Mejor calidad que coresets del mismo tamaño \\
\bottomrule
\end{tabularx}
\end{table}

Popcorn y Voevodski muestran los dos límites que el paralelismo no resuelve
(Tabla~\ref{tab:limites}). Popcorn conserva el algoritmo exacto y paraleliza casi todo su trabajo en
la GPU, así que la ley de \textcite{amdahl1967} no lo frena; lo frena la memoria, porque la
matriz kernel crece con $n^2$. Voevodski esquiva ese muro resumiendo los datos en unas pocas
semillas, pero paga con una garantía que depende de la estabilidad de los datos y con una parte
serial, $O(t^3)$, que acota el speedup; lo que ofrece a cambio es escalamiento débil, más datos en el
mismo tiempo \parencite{gustafson1988}, que proyecta sin medir en paralelo. Uno mantiene el algoritmo
y choca con la memoria; el otro cambia el algoritmo para no chocar y choca con Amdahl. Nuestro Lloyd exacto no enfrenta el muro de memoria de Popcorn ni necesita la aproximación de Voevodski: su límite será la fracción serial y la
granularidad del reparto (sección~\ref{sec:diseno}).

\begin{table}[H]
\caption{Qué limita a cada enfoque frente a nuestro Lloyd}\label{tab:limites}
\small
\begin{tabularx}{\textwidth}{@{}L{3cm}YYY@{}}
\toprule
 & \textbf{Popcorn} & \textbf{Voevodski} & \textbf{Nuestro Lloyd} \\
\midrule
Resultado & Exacto (Kernel K-means) & Aproximado: hasta 10 veces el óptimo en $k$-means, si hay
  estabilidad & Exacto \\
\addlinespace
Qué se paraleliza & Casi todo: matriz kernel y distancias en la GPU & Solo la pasada $O(n)$ de
  Voronoi & La asignación $O(nk)$ de cada iteración \\
\addlinespace
Qué lo limita & Memoria $O(n^2)$ de la matriz kernel & Parte serial $O(t^3)$ y supuesto de
  estabilidad & Fracción serial (lectura, barrera, reducción) y granularidad \\
\addlinespace
Memoria con 2,8 millones de viajes & Unos 32~TB & Resúmenes de $t$ semillas & Unos 136~MB \\
\bottomrule
\end{tabularx}

{\footnotesize\textit{Nota.} La memoria de Popcorn corresponde a la matriz kernel densa en precisión
simple; la de Lloyd, a las 6 features en doble precisión. Cálculo propio a partir de los tamaños
del dataset limpio.\par}
\end{table}

% --------------------------------------------------------
\subsection{Antecedentes y fuentes de apoyo}\label{sec:antecedentes}
% --------------------------------------------------------

Tres trabajos anteriores a la ventana de cinco años explican de dónde viene esta línea y se citan
solo como fundamento. k-means$\|$ \parencite{bahmani2012} paraleliza la inicialización en pocas
rondas y es la línea base de SOCCER. Yinyang K-means \parencite{ding2015} evita distancias con cotas
y garantiza el mismo resultado que K-means estándar, un antecedente directo de la poda de ASB. knor
\parencite{mhembere2017} adapta K-means a arquitecturas NUMA y minimiza las barreras de
sincronización, lo que plantea cuándo una sola máquina bien aprovechada compite con un clúster.

Entre las fuentes recientes que no forman parte de la cuota, \textcite{he2022} controlan el error de
redondeo al actualizar centroides con una suma en dos pasos, lo que respalda nuestra prueba de
equivalencia numérica. \textcite{mussabayev2023} discuten estrategias de K-means para big data,
\textcite{wan2021} aplican un K-means paralelo sobre Spark al diagnóstico de fallas y
\textcite{mclaughlin2023} proponen un clustering adaptativo paralelo para flujos de datos. La
documentación de NVIDIA \parencite{nvidia2025cuml,rapids2026dask} se usa solo como referencia de
implementaciones de producción, y sus cifras se leen como benchmarks del propio proveedor.
